% Ch2_4, slide 9: the dipole +q, -q (separation d) and a distant point P; R+, R, R-
\begin{tikzpicture}[scale=1]
  \coordinate (O) at (0,0); \coordinate (Qp) at (0,0.35); \coordinate (Qm) at (0,-0.35); \coordinate (P) at (6.2,2.0);
  \draw[ELax] (0,-0.9)--(0,2.6) node[ELlab,anchor=south]{$z$};
  \draw[ELthin] (Qp)--(P) (O)--(P) (Qm)--(P);
  \draw[ELvec,line width=1.2pt] (Qm)--(Qp);
  \node[ELlabs,anchor=east] at (-0.05,0) {$\vect{d}$};
  \fill[ELcSource] (Qp) circle (1.8pt); \fill[ELcField] (Qm) circle (1.8pt);
  \node[ELlabs,anchor=east] at (-0.1,0.42) {$+q$}; \node[ELlabs,anchor=east] at (-0.1,-0.42) {$-q$};
  \fill[ELink] (P) circle (1.6pt); \node[ELlab,anchor=west] at (P) {$P$};
  \node[ELlabs,anchor=south] at (2.4,1.15) {$R_+$};
  \node[ELlabs,anchor=north] at (2.6,0.8) {$R$};
  \node[ELlabs,anchor=north] at (2.2,0.35) {$R_-$};
  \draw[ELthin] (0,0.9) arc[start angle=90,end angle=17.9,radius=0.9];
  \node[ELlabs,anchor=south west] at (0.25,0.8) {$\theta$};
\end{tikzpicture}
